\subsection{A paired-block circuit for the bit correction}
\label{sec:paired-bit-construction}
Use $h=50$, $v=\binom h3=19600$, $N=v^3$, and $m=h^3=125000$.
The complex network above is retained at this ground size. For the bit
network, we replace the intersection-one side correction by a smaller circuit.
Fix a common point $i$. On the other $n=h-1=49$ points, inputs $x_{ab}$
represent source triples $\{i,a,b\}$. The required partial output is
\[
 z_{cd}=\sum_{\{a,b\}\cap\{c,d\}=\varnothing}x_{ab}.
\]
All scalar additions are in $\mathbb F_2$, but the computation below always
adds disjoint formal input supports.

It is useful to solve a weighted version. Give a graph an edge value $e_{uv}$
and a vertex value $w_u$, with pairwise disjoint supports. Return its total
value and the values left after deleting each one or two vertices, including
all their incident edges. At the root, take $w_u=0$ and $e_{uv}=x_{uv}$.
Group the ordered vertices consecutively into pairs, with a final singleton
when necessary. A block $A$ becomes a vertex of a smaller graph with weight
\[
 \widetilde w_A=\sum_{u\in A}w_u+\sum_{\{u,v\}\subset A}e_{uv},
 \qquad
 \widetilde e_{AB}=\sum_{u\in A,\ v\in B}e_{uv}.
\]
Recursively compute the smaller graph's total, its value $O_A$ after deleting
block $A$, and its value $F_{AB}$ after deleting both $A,B$.

For $a\in A$, put
\[
 C_a=\sum_{u\in A\setminus\{a\}}w_u,\qquad
 E_{a,B}=\sum_{u\in A\setminus\{a\},\ v\in B}e_{uv},\qquad
 H_{a;B}=C_a+\sum_{C\ne A,B}E_{a,C}.
\]
Compute all $H_{a;B}$ by the leave-one-out prefix/suffix routine on the
vector $(C_a,(E_{a,C})_{C\ne A})$, requesting only the outputs that omit
one of the $E$ terms. The total of this vector is denoted $H_a$.
If $a,b$ are in different blocks $A,B$, their exclusion output is
\[
 F_{AB}+H_{a;B}+H_{b;A}
   +\sum_{u\in A\setminus\{a\},\ v\in B\setminus\{b\}}e_{uv}.
\]
For two vertices in one block it is $O_A$; the single exclusion of $a$ is
$O_A+H_a$. These are disjoint partitions of the surviving weighted graph.
This proves the recursion by induction. The weights created by internal edges
remain inside the recursion, avoiding a separate exclusion circuit for them.

Use direct sums for at most four vertices. Compute other totals by balanced
binary addition. In the four-term formula, first add $F_{AB}+H_{a;B}$,
then add the other two terms, then combine those two results. Intern equal
formal supports and remove unused nodes. At $n=49$, the exact circuit has
9813 additions and 1176 designated outputs. All 1271256 nonzero coefficients,
and every required zero, are checked by comparing all output supports with
their definitions.

Identify equal inputs and equal sums across the $h$ common-point groups.
A sum shared by groups $i\ne j$ consists entirely of triples $\{i,j,k\}$;
its fixed pair and variable vertices identify it exactly. Keep the first
encountered decomposition, in common-point and then local-node order. Retain
all $q=h\binom{h-1}{2}=58800$ partial outputs, injecting their sum into each
physical target. Since every neighbor pair has a unique common point, this
remains exactly the intersection-one matrix. The global graph has
\[
 c=450394,\qquad q=58800,\qquad c+q=509194.
\]
Here 40256 additions have been merged from the separate local graphs.

\subsection{Reversible roles and frames in both directions}
For every node, regard its later uses and designated outputs as outgoing edges.
At an input node, allocate one role per outgoing edge and copy a pivot into
the others. At an addition node, add the second incoming role into the first,
reuse the first as one outgoing pivot, and copy it into fresh roles for the
other outgoing edges. The other input role retires. These are invertible
pointwise gates. There are $2c+q$ outgoing uses and $c$ pivot identifications,
so $R=c+q=509194$ roles suffice. Roles are compiled from the merged graph;
no simultaneously occupied physical slots are identified.

Let $L$ be the product of these gates, $V$ the source copy, $J$ the injection
of the designated partial outputs, and $G,R_0$ the original central operations.
The chronological schedule is
\[
 L,\ J,\ L^{-1},\ R_0,\ V,\ G,\ R_0,\ L,\ J,\ L^{-1},\ G,\ V.
\]
The two injections give $JLz+JL(z+Vx)=JLVx$ for arbitrary initial scratch $z$.
The cleanup restores $z$ and all central scratch. Since $JLV+R_0G=I$,
the invocation is $y\gets y+x$. Reverse and invert the schedule for stage 2.
Forward, inverse, forward invocations therefore exchange the two data banks.

Use the original rational form $I-J/9$ on $F=\mathbb Q^{50}$. At a stage,
write $A=B\mathbin\perp P$ as before, suppressing the future line factor $Q$.
Set $D_0=B\otimes F$, $D_1=A\otimes F$, and
$D_U=(B\otimes F)\mathbin\perp(P\otimes U)$.
Let $U_z$ be the span of the source indicators reaching node $z$.
Every retained node is a sum from an original common-point group, so all its
triples contain some $i$. For a vector $u$ in that span,
\[
 \sum_j u_j=3u_i,\qquad
 \langle u,u\rangle=\sum_{j\ne i}u_j^2>0\quad(u\ne0).
\]
Thus $U_z$ is nondegenerate. Along every graph edge $z\to w$,
$U_z\subset U_w$. Forward middle gates use $D_{U_z}$; the physical roles
follow paths, so these frames increase. A designated output for target $T$
has $U_z\subset t_T^\perp$, because every contributing triple is a neighbor.

For the reversed middle computation use $D_{U_z^\perp}$. The ambient form
is nondegenerate, so all these complements are nondegenerate. Reversing an
edge reverses inclusion: $U_w^\perp\subset U_z^\perp$.
The original output begins at the physical $X$ line
$\langle t_T\rangle\subset U_z^\perp$; an original input ends at exactly
its physical $Y$ complement $t_S^\perp$. No positivity assumption on the
reachable-target spans is needed. At every gate all incident roles receive
the same frame. Early mixers use $D_0$, and late cleanup uses $D_1$.
Newly active roles enter from $D_0$ and retired roles grow to $D_1$.
Tensoring by $Q$ gives the original data-stage boundaries in both orientations.
The only decreases remain the central returns, each of dimension $h$.

Share the complete auxiliary bank of a stage-1 invocation $(A,B)$ with stage 3
at $(B,\pi(A))$, matching local role indices. Here $\pi$ is a bijection of
triples with $|T\cap\pi(T)|=1$: pair consecutive ground points; for a triple
with a full pair and singleton, keep the singleton and cycle the full pair
among pairs not containing it; otherwise keep the point in the least indexed
pair and flip the other two to their partners. Each rule is bijective on its
class. All invocations restore arbitrary scratch, and the joining labels are
\[
 E=F\otimes\langle t_A\rangle\otimes\langle t_B\rangle,
 \qquad H=\langle t_B\otimes t_{\pi(A)}\rangle^\perp\otimes F.
\]
They are nondegenerate and $E\subset H$ by the chosen orthogonality.
Each join removes one auxiliary role and exactly $m$ of edge rank.

\begin{proposition}[Paired-block bit interface]\label{prop:paired-bit-interface}
The new bit network restores arbitrary scratch, exchanges the data banks,
and satisfies the original all-role endpoint identity with
\begin{align*}
 W_{\rm b}^{\rm pair}&=2N+2v^2(R+h)=406321422080000,\\
 s_{\rm b}^{\rm pair}&=W_{\rm b}^{\rm pair}m-N+6v^2h^2
                         =50790175992864000000,\\
 W_{\rm b}^{\rm pair}m-s_{\rm b}^{\rm pair}&=1767136000000>0.
\end{align*}
\end{proposition}
\begin{proof}
The scalar schedule and both frame assignments were verified above. The total
decreasing dimension is $3v^2h^2$. Telescoping gives total absolute label change
$W_{\rm b}^{\rm pair}m-2N+6v^2h^2$; the unchanged negative source projections
add $N$ rational rank. Data endpoints are unchanged and every surviving
auxiliary role has endpoints zero and $I_m$, with its scalar value restored.
The original common-frame identity and rational-shear transfer apply. All
circuits, rational bases, and gates are fixed finite data.
\end{proof}

\subsection{Explicit exponents at ground size fifty}
Take $a=296/10^{11}$, $a_{\rm c}=1/10^{11}$, and $L_0=11737/1000$.
The exact deficits satisfy
\[
 \eta_{\rm b}^{\rm pair}=\frac{23}{661055000}>aL_0,\qquad
 \eta_{\rm c}=\frac{239}{1202215191250}>a_{\rm c}L_0.
\]
For $1\le x\le2$, put $z=(x-1)/(x+1)$,
$S(x)=2\sum_{j=0}^{23}z^{2j+1}/(2j+1)$, and
$E_0(x)=2z^{49}/(49(1-z^2))$. Then
$S(x)\le\log x\le S(x)+E_0(x)$, and exact rational comparison gives
\[
 \log m\le16(S(2)+E_0(2))+S(m/2^{16})+E_0(m/2^{16})<L_0.
\]
The inequality $e^{-x}>1-x$ therefore proves the required strict bounds with
\begin{equation}\label{eq:explicit-motif-exponents}
 \tau=1-\frac{296}{10^{11}},\qquad \sigma=1-\frac1{10^{11}}.
\end{equation}
The complex network itself is unchanged apart from the ground size $h=50$;
its construction, binary residual bases, and scalar corrections remain those
proved above.
